\chapter{Rebuilding VidyaRAG Yourself}
\label{ch:rebuilding}

\begin{recommended}[how to use this chapter]
This is a suggested order for rebuilding the system from an empty folder, as a way to understand it deeply. The stages follow the order in which the repository was actually built (Chapter~\ref{ch:history}). Code blocks marked \emph{illustrative sketch} are simplified teaching examples, \textbf{not} the repository's code; the real implementations are referenced beside them. Each stage ends with an acceptance check and the pitfalls this project hit.
\end{recommended}

\section*{Stage map}

\begin{center}
\begin{tabularx}{\textwidth}{@{}c>{\raggedright\arraybackslash}p{4.4cm}L@{}}
\toprule
\textbf{Stage} & \textbf{Goal} & \textbf{Real code to study}\\
\midrule
0 & Project skeleton & \code{pyproject.toml}, \code{settings.py}, \code{config/}\\
1 & Fetch the corpus & \code{ingest/corpus.py}, \code{ingest/download.py}\\
2 & Parse with structure & \code{ingest/parse.py}\\
3 & Chunk & \code{ingest/chunk.py}, \code{llm/provider.py}\\
4 & Embed and store & \code{store/}, \code{ingest/pipeline.py}\\
5 & Baseline RAG with citations and tracing & \code{retrieve/dense.py}, \code{generate/}, \code{observe/trace.py}, \code{pipeline.py}\\
6 & Evaluation harness & \code{evaluation/}\\
7 & Reranking ablation & \code{retrieve/rerank.py}, \code{config/profiles/rerank.yaml}\\
8 & A negative experiment: decomposition & \code{retrieve/decompose.py}\\
9 & Self-check loop & \code{correct/}\\
10 & Guardrails & \code{guard/}\\
11 & Interfaces & \code{api/}, \code{cli.py}, \code{app/app.py}\\
12 & Packaging and deployment & \code{Dockerfile}, \code{.github/workflows/}, \code{scripts/deploy_space.py}\\
\bottomrule
\end{tabularx}
\end{center}

\section{Stage 0 --- skeleton}

Create the project with uv, a \code{src/} layout and two configuration layers: environment settings for secrets and endpoints, YAML profiles for behaviour, validated with \code{extra="forbid"}. Add ruff, black, mypy, pytest and gitleaks from day one.

\textbf{Acceptance:} \code{uv run pytest} passes on an empty test; a misspelt profile key fails at load. \textbf{Pitfall:} decide the oldest Python you must support before writing code --- the Space's Python 3.10 forced a compatibility shim late.

\section{Stage 1 --- corpus}

Declare the books once (slug, edition, URLs, licence). Download with streaming, a \code{.part} file renamed on completion, HTTP \code{Range} resumption and a SHA-256 manifest.

\textbf{Acceptance:} interrupting and re-running resumes; a second run makes no request. \textbf{Pitfalls:} verify the licence \emph{per edition} with the publisher's API; handle 416 and ignored \code{Range} headers.

\section{Stage 2 --- parsing with structure}

\begin{lstlisting}[language=Python]
# Illustrative sketch -- not the repository's code (see ingest/parse.py)
toc = sorted((p, lvl, title) for lvl, title, p in doc.get_toc(simple=True) if p >= 1)
for page in doc:
    blocks = [b for b in page.get_text("blocks") if b[6] == 0 and inside_body_band(b, page)]
    text = join_in_reading_order(blocks)          # detect a central gutter for two columns
    section = section_for(page.number + 1, toc)   # most recent outline entry at or before
\end{lstlisting}

\textbf{Acceptance:} every page carries a chapter, section and printed page; answer keys and question sets are gone. \textbf{Pitfalls:} sort the outline (a stray entry can swallow the book); normalise non-breaking spaces; read printed page numbers from running heads; make exclusion inherit downwards, not sideways.

\section{Stage 3 --- chunking}

Split into sentences, then pack whole sentences into chunks within a section, measured with the \emph{embedding model's} tokenizer, with a sentence-level overlap.

\textbf{Acceptance:} no chunk exceeds the model's limit; ids are stable across runs. \textbf{Pitfalls:} a different tokenizer silently truncates chunks; the tokenizer's special tokens count too; abbreviations need re-joining; a regular expression look-behind must be fixed-width.

\section{Stage 4 --- embedding and storage}

\begin{lstlisting}[language=Python]
# Illustrative sketch -- not the repository's code (see store/collection.py)
point_id = str(uuid.uuid5(NAMESPACE, chunk.chunk_id))     # re-running overwrites, never duplicates
client.upsert(collection, [PointStruct(id=point_id, vector={"dense": vec}, payload=citation_payload)])
\end{lstlisting}

\textbf{Acceptance:} a second ingest writes nothing new; a dimension mismatch fails with a clear message. \textbf{Pitfalls:} carry the licence and citation on every point; write run reports beside the index they describe.

\section{Stage 5 --- baseline RAG}

\begin{lstlisting}[language=Python]
# Illustrative sketch -- not the repository's code (see pipeline.py, generate/)
hits = retrieve_dense(question, limit=20)[:5]
prompt = RULES + "\n\n".join(f"[{i}] {h.citation}\n{h.text}" for i, h in enumerate(hits, 1))
raw = llm.generate(system=RULES, user=prompt + f"\n\nQuestion: {question}", temperature=0)
citations = [hits[n - 1] for n in markers(raw) if 1 <= n <= len(hits)]   # drop invented markers
\end{lstlisting}

Record every stage's time and every model call's tokens in a trace returned with the answer.

\textbf{Acceptance:} every displayed citation points at a supplied passage; empty retrieval never calls the model. \textbf{Pitfalls:} parse grouped markers like \code{[1, 2]}; \textbf{do not nest timed stages without excluding them from the total}; record usage for every model call, not only generation; version the prompt.

\section{Stage 6 --- evaluation harness}

Build the gold set with provenance and validators; compute recall@k, hit rate, MRR and recall@context from ids; add RAGAS; add abstention precision, recall and false-abstention rate; cache answers; withhold metrics from runs with too many failures.

\begin{lstlisting}[language=Python]
# Illustrative sketch of a judge that cannot silently fail -- not the repository's code
reply = client.chat.completions.create(model=m, messages=msgs, temperature=0, max_tokens=16)
label = (reply.choices[0].message.content or "").strip().strip("*").upper()
if not label.startswith(("REF", "ANS")):
    raise JudgeError(f"unparseable label: {label!r}")     # count it, do not guess
return label.startswith("REF")
\end{lstlisting}

\textbf{Acceptance:} a system that refuses everything scores badly; a test feeds the judge truncated and decorated labels; the judge's agreement with a small hand-labelled set is reported. \textbf{Pitfalls:} score each metric from the right list; keep the cache key complete; never quote partial runs; commit every run a document cites.

\section{Stage 7 --- reranking, as an ablation}

Add one flag and one profile; predict that pool recall will not move; run the gold set; read results by question type.

\textbf{Acceptance:} pool recall unchanged, recall@context and MRR up. \textbf{Pitfall:} the mean latency will be dominated by model loading unless a warm-up question is excluded.

\section{Stage 8 --- a negative result}

Build decomposition with reciprocal rank fusion, measure it, and be willing to reject it. Diagnose \emph{why} it failed before trying a fix.

\section{Stage 9 --- the self-check}

Grade claims against passages with a different model; decide with explicit thresholds and a bounded retry that re-retrieves; detect refusals explicitly.

\textbf{Acceptance:} tests for accept, retry, abstain, refusal and grader failure. \textbf{Pitfalls:} a refusal scores as perfectly grounded; decide deliberately what an \emph{unverified} answer looks like to the user; before crediting the loop with refusals, check whether the baseline already refuses.

\section{Stage 10 --- guardrails}

Screen questions before retrieval and passages before the prompt; measure false positives on the whole corpus before worrying about recall.

\textbf{Acceptance:} zero corpus false positives, and a written list of known bypasses. \textbf{Pitfall:} do not describe a handful of test cases as a measured block rate.

\section{Stage 11 --- interfaces}

One pipeline per process (the embedded index locks); an API with typed request and response models; a CLI for workflows; a demo that shows the working.

\textbf{Pitfalls:} keep one source of truth for configuration in every entry point; make the response say whether an answer was verified, abstained or blocked.

\section{Stage 12 --- packaging and deployment}

Multi-stage Docker build with the environment built where it runs; health check over HTTP; CI that actually starts the container; a deployment script that never writes secrets to files.

\textbf{Pitfalls (all hit by this project):} console-script shebangs, editable installs, configuration paths inside site-packages, read-only mounts for a store that needs a lock file, Python version pins on the hosting platform, and lock files copied into deployments.
